% Plan de réflexion LaTeX
\documentclass[12pt]{article}
\usepackage[utf8]{inputenc}
\usepackage[T1]{fontenc}
\usepackage[french]{babel}
\begin{document}

\section*{Introduction}
Objectifs lors de l'introduction :
\begin{itemize}
    \item Capteur ; d'attention generalement
    \item Exposer et Definir la problematique
\end{itemize}
\\
Les médias occupent une place centrale dans nos sociétés contemporaines. Avec l'essor d'Internet et des réseaux sociaux, le nombre de sources d'information s'est considérablement multiplié. Cette diversité permet d'accéder à une multitude de points de vue, mais elle favorise aussi l'apparition de croyances variées, parfois contradictoires.


\subsection*{\textbf{Problématique :} Dans quelle mesure la diversité des médias facilite-t-elle l’accès à l’information fiable pour les citoyens ?}

\section*{Plan de réflexion}

\begin{enumerate}
    \item \textbf{I. Diversification des sources d'information}
    \begin{itemize}
        \item Accès facilité à l'information
        \item Pluralité des points de vue
    \end{itemize}
    \item \textbf{II. Multiplication des formes de croyance}
    \begin{itemize}
        \item Apparition de croyances alternatives
        \item Influence des réseaux sociaux et des fake news
    \end{itemize}
    \item \textbf{III. Conséquences sur la société}
    \begin{itemize}
        \item Fragmentation des opinions
        \item Difficulté à distinguer le vrai du faux
    \end{itemize}
\end{enumerate}
\section*{Conclusion}
Objectifs de la conclusion :
\begin{itemize}
    \item Amener de nouvelles perspectives sur la problématique.

    En conclusion, la multiplication des sources d'information a conduit à une diversification des formes de croyance. 
    Cette pluralité enrichit le débat public, mais pose aussi des défis pour la cohésion sociale et la recherche de la vérité. 
    Il est donc essentiel de développer l'esprit critique pour mieux naviguer dans cet univers médiatique.

\end{document}

